\subsection{Zukünftige Entwicklungen und Datennutzung}

Aus den erzielten Ergebnissen ergeben sich wesentliche Möglichkeiten für zukünftige Arbeiten. Der primäre nächste Schritt besteht in der flächendeckenden Anwendung des Workflows auf das gesamte Stadtgebiet von Leipzig. Diese großflächigen Datensätze dienen nicht nur der Vervollständigung des Datenbestands, sondern sollen im nächsten Schritt aktiv genutzt werden, um forstwirtschaftliche und stadtplanerische Entscheidungen datenbasiert zu unterstützen. Das größte Potenzial dieser Datennutzung liegt in der Anwendung der fusionierten geometrischen und botanischen Attribute für maschinelle Lernverfahren. Der erstellte Datensatz aus Geometrie- und Katasterdaten eignet sich als Trainingsgrundlage für Klassifikationsmodelle. Ziel ist eine automatisierte Bestimmung der Baumgattung basierend auf LiDAR-Rohdaten. Eine solche algorithmische Artenerkennung würde die kontinuierliche Aktualisierung des städtischen Baumkatasters, insbesondere in unzugänglichen Bereichen erleichtern. Darüber hinaus ermöglicht die Einbindung multi-temporaler LiDAR-Daten eine dynamische Vegetationsanalyse. Durch den Abgleich mehrjähriger Punktwolken lassen sich Wachstum und Vitalitätsveränderungen im urbanen Raum automatisiert quantifizieren. Dies erlaubt es, Bewässerungs- und Pflegemaßnahmen gezielt zu optimieren sowie bei beginnenden Vitalitätseinbrüchen frühzeitig Erhaltungsmaßnahmen einzuleiten. 
